\documentclass[11pt]{article}
\usepackage[a4paper,margin=25mm]{geometry}
\usepackage{amsmath,amssymb,amsthm,booktabs}
\newtheorem{theorem}{Theorem}
\newtheorem{lemma}{Lemma}
\newcommand{\eps}{\varepsilon}
\newcommand{\sigs}{\sigma}
\begin{document}
\title{Finite source bounds beyond $z=1$ and the effect of retained squares}
\author{}\date{23 October 2026}\maketitle

This note concerns the local C1 certificate in the recovered order669
construction. Its equations (P1), (SP3), (SP6), and the positive twelve-root
and six-root quadratic forms are used with their literal definitions.
The conclusion is source admission to those local certificates. It is
not an inclusion theorem for the full quantum correlation set.

\begin{theorem}[A common finite-$R$ box]\label{thm:R}
Fix $R\ge1$, put $T=1+R$, and set
\[
 \eps_R=2^{-96}T^{-16},\qquad
 \delta_R=2^{-28}T^{-4},\qquad M_R=2^{16}T^4.
\]
Suppose that (P1) is evaluated on
\[
 0<x,\sigs\le\eps_R,\quad 0<z\le R,\quad
 0<b,h<1,\quad r=bz<1,\quad E,M,J>0,\quad D,\beta\ge0.
\]
There is no bound on $H=h/\sigs$.
For $\lambda>0$ the complete twelve-root Gram has
\[
 Q(v,w)=Q_0+\sum_{\mathcal A}C_{ij}v^iw^j,
 \quad v=\sqrt u,\quad
 \mathcal A=\{(0,1),(0,2),(1,0),(2,0),(2,1),(4,0),(4,1)\},
\]
with reference $B\succeq I/3$ and
\begin{equation}\label{eq:mainbounds}
 \|Q_0-B\|_\infty\le M_R\sqrt{\eps_R},\qquad
 \|C_{ij}\|_\infty\le M_R.
\end{equation}
Consequently its separated testing construction is valid on
$0\le\sqrt u,w\le\delta_R$.
For $\lambda\le0$ the six-root construction satisfies all its source
guards, with the same testing box. In particular $d_G\ge1$ and
$4\kappa H\ge1$. The assertion is literal at $\beta=0$ and, on the
positive branch, at $D=0$. On the nonpositive branch $D=0$ is excluded by
$d_G\ge1$, rather than by cancellation of $D$.
\end{theorem}

All constants above are real constants shared by the source and testing
fields; no source parameter is identified with a testing coordinate.
For rational data one may first replace $R$ by an integer majorant.
The finite testing list and its degrees do not increase with $R$;
its testing interval changes.

\paragraph{Source estimates.}
Write $\eps=\eps_R$, $Q_b=1-b$, $\theta=1-h$, $k=\sqrt x$,
$q=\sqrt z$, $e=xh/\sigs$, and $H_0=1+zQ_b$. The following deliberately
loose bounds imply the theorem:
\begin{align}
 &1\le P\le1+4T\eps,\quad |J-1|\le\eps,\quad |M-1|\le T\eps,
 \quad 1\le B_s\le1+6T\eps,\quad 1\le H_0\le T,\label{eq:src1}\\
 &e\le5T\eps,\quad |\mathcal H-H_0|\le32T^2\eps,
 \quad D\le3T\sigs^2,\quad \beta\theta\le h,\label{eq:src2}\\
 &\tfrac13\le p\le T,\quad \tfrac13\le l,\kappa\le1,
 \quad pb\le2,\quad C(1-r)\le T,\quad m\le5T^2\eps^2,\label{eq:src3}\\
 &|\nu|\le2T,\quad |\lambda_0|\le4T,\quad
 |\lambda_0+2\theta H_0|\le64T^2\eps,\label{eq:src4}\\
 &|K_A-(2\theta H_0-2h(1-r))|\le1024T^4\eps,\label{eq:src5}\\
 &|R_0|\le64T(x+\sigs),\quad |B_x|\le40T\eps,\quad
 |A_y|,|\chi|\le64T^2,\quad |R_x|\le24T,\quad |k_1|\le16T^2.
 \label{eq:src6}
\end{align}
Here and below an expression such as $C(1-r)$ is bounded before taking
absolute values, so no lower bound on $1-r$ is assumed.

For completeness, these estimates follow from the native identities by
the following inequalities. The positive expansion of $P$ gives
$P-1\le2T(x+\sigs)$, using $bz=r<1$ in the term
$\sigs bx^2hz^2$. Also $J-\sigs>0$ and
$|MP-1|\le6T\eps$, so $(MP)^{-1}\le2$ and
$|(MP)^{-1}-1|\le12T\eps$. Initially $e<1$ and
$\mathcal H\le4T$. The two exact identities
\[
 2\theta P\beta=hB_s-\sigs H_0,\qquad
 2P(JD+x\theta\beta)=\sigs^2\mathcal H
\]
give
\[
 eB_s\le xH_0+\sigs\mathcal H\le Tx+4T\sigs.
\]
Expansion of $\mathcal H$ then yields
\[
 |\mathcal H-H_0|\le4Tx+3Te+2Te^2\le32T^2\eps.
\]
The second identity gives $D\le3T\sigs^2$, and the first gives
$\beta\theta\le hB_s/(2P)\le h$. The numerator of $p$ is at least
$H_0(1-e)$, and its product by $b$ is at most $1+2\eps$.
The numerator of $l$ is bounded below by $M(1-e)$ and above by
$(1+\sigs r)\max\{M,(1+xr)(1+\sigs z)\}$.
These prove \eqref{eq:src3}; in particular the bounds for $p$ and $pb$
are used separately.

The pole-free formulas
\begin{align*}
 \nu&=2D/\sigs+\{H_0-ez(J-\sigs)\}/P,\\
 \lambda_0&=2e\theta z(J-\sigs)/P-2\theta H_0/P
                         +8Dh/\sigs-4bp\sigs z
\end{align*}
give \eqref{eq:src4}. Put
\[
 N_\kappa=(1+x)(1+\sigs r)(1+\sigs z)-MP=xF_x+\sigs F_\sigs,
 \quad F_\sigs=r-1+\sigs rz.
\]
The polynomial $F_x$ has coefficient absolute sum $37$ and degree at
most three in $z$. Thus $|F_x|\le37T^3$. To verify this finite bound
without any rational expansion, its coefficients in $x$ are
\begin{align*}
 F_x={}&b\sigs^2z^2-b\sigs z^2-b\sigs z-bz
 -h\sigs z^2-3h\sigs z-hz+\sigs z^2+2\sigs z+z\\
 &+x\{-b^2\sigs z^2-4bh\sigs z^2+b\sigs z^2-bz
 -2h^2\sigs z^2+h^2z+2h\sigs z^2-2hz+z\}\\
 &+x^2\{-b^2h\sigs z^3-bh^2\sigs z^3+bh^2z^2
 +bh\sigs z^3-bhz^2+h^3z^2-2h^2z^2+hz^2\}.
\end{align*}
The exact cancellation
$H(2\kappa-1)=(eF_x+hF_\sigs)/(MP)$ gives
\[
 |H(2\kappa-1)-h(r-1)|\le384T^4\eps.
\]
Together with $K_A=2H(2\kappa-1)-\lambda_0$ this proves
\eqref{eq:src5} for arbitrary $H$.
Finally use the exact cancellation
\begin{align*}
 -PR_0/4={}&z(1-e)\{1+xh-(1+x)(1+\sigs r)\}\\
 &+(1-r)\{1+xhz-P\}-2DP(1-r).
\end{align*}
The three terms are bounded respectively by
$T(x+2\sigs)$, $3Tx+2T\sigs$, and $T\sigs$.
This proves the $R_0$ bound. Substitution in (SP3), with
$|h+b-1|\le1$, proves the remaining \eqref{eq:src6}; no negative
power of $z$, $b$, $Q_b$, or $\theta$ is used in those estimates.

\paragraph{Positive branch and angular coefficients.}
Put $d_G=4D/\sigs^2$, $\rho=\lambda x/\sigs$,
$S=d_G+\rho$, $L_0=\lambda\sigs$, $q_0=\lambda k$, and
$c=\max(0,1-K_A)$. The exact source relation
\begin{align*}
 S={}&2\mathcal H/P+d_G(1-J+2hx)+2eH_0/P\\
 &-2e^2z(J-\sigs)/P-4bpxz
\end{align*}
gives
\[
 |S-2H_0|\le128T^2\eps,\quad H_0\le S\le3T,
 \quad L_0\le3,\quad q_0^2=L_0\rho\le9T.
\]
Indeed the five respective errors from $2H_0$ are at most
$72T^2\eps$, $36T\eps$, $10T^2\eps$,
$T^2\eps$, and $8T\eps$.
The ratios $\sqrt{d_G}/S$, $\sqrt\rho/S$, $d_G/S$, and
$\rho/S$ are at most one.
Also
\[
 (x+\sigs)^2H=xe+2\sigs e+\sigs h\le2\eps,
 \qquad |R_0\sqrt\lambda|\le256T\sqrt\eps.
\]
Although $a_0=k/\sigs$ need not be bounded, the angular factors obey
\begin{equation}\label{eq:angular}
 0\le c\le4,\quad
 \frac{ca_0}{S}\le32,\quad
 \frac{c}{SL_0}\le32,\quad
 \frac c\lambda\le32\sigs.
\end{equation}
If $c>0$, \eqref{eq:src5} forces $h>1/8$; otherwise its reference
value is at least $3/2$. Then $L_0\ge h>1/8$ and
\[
 a_0^2(hB_s-\sigs H_0)\le\mathcal H,\quad
 hB_s-\sigs H_0\ge1/16,\quad
 a_0^2\le32H_0.
\]
Since $S\ge H_0\ge1$, these imply \eqref{eq:angular}.
If $c=0$, each coefficient in \eqref{eq:angular} is exactly zero.

\paragraph{A finite bound on the complete twelve-root correction.}
Here is an explicit coefficient accounting, so the powers of $T$ do
not rest on a continuity argument. Treat the twelve root labels
\[
 (N_x,N_y,R,\Gamma_*,A_\theta,T_\theta,d_1,d_2,d_3,d_4,U_*,K_*)
\]
as formal independent variables, solely to extract Gram coefficients.
Let $\mathcal N$ be the sum of absolute coefficients over both the
testing monomials in $v,w$ and quadratic root monomials. Let
$\mathcal N_0$ retain only testing degree zero. The symmetric-matrix
row norm of every coefficient is at most $\mathcal N$.
The exact conductor expressions are
\begin{align*}
 F_1={}&\frac{\sqrt{d_G}}S u\Gamma_*
 +\frac{\sqrt\rho}S\{k d_1-u d_2+(u-k)d_3\}
 -\frac\rho S\sqrt\sigs T_\theta,\\
 F_2={}&\frac{R_0\sqrt{d_G}}S qw\Gamma_*
 +\frac{R_0\rho\sigs z}S R
 -\frac{R_0\sqrt\rho}S(qw+z\sqrt\sigs)(d_2-d_3)\\
 &+\left(\frac{R_0q_0z}S+
                \frac{R_0\sqrt\lambda\sqrt\rho}S qw\right)T_\theta
 -8pquwF_1 .
\end{align*}
They satisfy
\[
 \mathcal N(F_1)\le4,\quad \mathcal N_0(F_1)\le3\sqrt\eps,
 \quad \mathcal N(F_2)\le128T^3,\quad
 \mathcal N_0(F_2)\le1024T^3\eps.
\]
Put $J_C=N_y+bR$ and $a_C=cd_G/S$. The full correction is
\begin{align*}
 \mathcal E={}&E_{yy}N_y^2+E_{RR}R^2+E_{yR}N_yR
 -4\nu N_xF_1+J_CF_2\\
 &-4uqw(2m+2D+1)N_xN_y+uqwR_xN_xR+qwB_xN_xA_\theta\\
 &+uA_yN_yA_\theta+\chi uA_\theta R
 -\tfrac12\rho\sigs R^2+\sqrt{q_0}\sqrt u\,U_*R\\
 &-a_C(2/\lambda+u+k)\sum_{i=1}^4d_i^2-a_Cu^2R^2\\
 &-\frac cS\{2a_0u+(u+k)/L_0\}\Gamma_*^2
 -\frac cS\{\sqrt{a_0u}\,U_*+\sqrt{a_0k}\,K_*\}^2.
\end{align*}
The reference $B$ is precisely the one in the order669 construction:
$N_x$ entry two; radial block
$\bigl(\begin{smallmatrix}2&2b-1\\2b-1&b^2+Q_b^2\end{smallmatrix}\bigr)$;
$\Gamma_*$ entry one; angular entries $(K_A+c)/2$; four entries $1/2$;
and final entries $1,1/2$. The radial determinant is one and trace at
most three, proving $B\succeq I/3$.
The definitions of $E_{yy},E_{RR},E_{yR}$ in (SP6) give the following
complete bounds; each row follows by multiplication of the preceding
displayed inequalities:
\[
\begin{array}{l|c|c}
 \text{part of }\mathcal E&\mathcal N&\mathcal N_0/\sqrt\eps\\\hline
 E_{yy}N_y^2+E_{RR}R^2+E_{yR}N_yR&64T^3&128T^3\\
 -4\nu N_xF_1&32T&24T\\
 J_CF_2&256T^3&2048T^3\\
 \text{three terms containing }qwN_x&64T^2&0\\
 uA_yN_yA_\theta+\chi uA_\theta R&128T^2&0\\
 -\rho\sigs R^2/2+\sqrt{q_0u}\,U_*R&3T&2T\\
 -a_C(2/\lambda+u+k)\sum d_i^2&32&272\\
 -a_Cu^2R^2&4&0\\
 -c\{2a_0u+(u+k)/L_0\}\Gamma_*^2/S&128&32\\
 -c\{\sqrt{a_0u}U_*+\sqrt{a_0k}K_*\}^2/S&128&32
\end{array}
\]
For example the constant radial coefficients obey
$|E_{RR,0}|,|E_{yR,0}|\le64T^3\eps$ after grouping
$C(1-r)$; their nonconstant coefficients have total at most $32T^3$.
The circuit row uses $2a_C/\lambda\le64\sigs$; the final two rows use
\eqref{eq:angular}. The sum of either column is less than
$2^{12}T^4<M_R$. This proves \eqref{eq:mainbounds}, including every
coefficient, rather than only a pointwise value of $Q$.

Set $r_*=M_R\sum_{\mathcal A}\delta_R^{i+j}$.
Then $r_*\le7M_R\delta_R=7/4096<1/12$ and
$M_R\sqrt\eps_R=2^{-32}T^{-4}<1/12$. Hence
\[
 Q_0-r_*I\succ I/6,\qquad
 \frac{\delta_R^{i+j}}2(M_RI\pm C_{ij})\succeq0.
\]
Multiplication by the testing squares
$1\pm(v/\delta_R)^i(w/\delta_R)^j$ gives the exact separated identity
for $Q(v,w)$.

\paragraph{The nonpositive branch.}
If $\lambda\le0$, its defining equation and $pb\le2$ imply
\[
 H^2\le2pbz\le4T,\quad \beta\theta^2\le2T\sigs^2,
 \quad \theta\ge\tfrac12,\quad -8T\le L=\lambda/\sigs\le0.
\]
Thus $s_L=\sqrt{1+L^2}\le9T$ and
$H\ge H_0/B_s\ge3/4$, so $4\kappa H\ge1$.
The second source identity gives the quantitative bound
\[
 d_G=\frac{2\mathcal H/P-4x\theta\beta/\sigs^2}{J}
       \ge 2H_0-128T^2\eps\ge1.
\]
Also $4a,4\ell\ge4/3$. In the six-root form use the roots
$(N_x,N_y,R,\Gamma,A_\theta,\widehat V)$, with
$\widehat V=s_Luw(u^2-x)W$.
Its reference has the same radial block, $N_x$ entry two,
$\Gamma$ entry $d_G$, $A_\theta$ entry $4\kappa H$,
$\widehat V$ entry one, and off-diagonal $L/(2s_L)$ between
$\Gamma$ and $\widehat V$. It is at least $I/3$.
For its seven corrections the coefficient bounds, using the complete
six-root list, are
\[
\begin{array}{c|rrrrrrrr}
 (i,j)&(0,0)&(1,0)&(0,1)&(1,1)&(2,0)&(0,2)&(2,1)&(3,0)\\\hline
 \mathcal N&256T^3\eps&256T^2&512T^2&32T^2&32T^3&11T&16T^2&81T^2
\end{array}
\]
The $(0,0)$ entries are $E_{RR,0}$, $E_{yR,0}/2$, and
$-s_Lx/2$; the latter contributes at most $9T\eps$ to
$\mathcal N$. For $(1,0)$ the three quadratic coefficients are
$-4\nu$, $A_y$, and $\chi+\lambda+s_L^2\sigs x$.
For $(2,0)$ they are the three radial coefficients, $L$,
$-s_L^2\sigs$, and $s_L$. These observations prove the largest
entries in the table directly; the others follow from
$q\le T$, $p\le T$, and \eqref{eq:src6}.
Every nonconstant row is below $M_R$, and the constant row below
$M_R\sqrt\eps_R$. At $0\le u,w\le\delta_R$ every normalized
nonconstant correction is at most $M_R\delta_R=1/4096<1/56$.
The six-root source guards and its separated construction therefore
hold. The positive branch's box has $u\le\delta_R^2\le\delta_R$,
so it is a common testing box for both branches.

\begin{lemma}[A limitation of discarding retained roots]\label{lem:retained}
No positive testing interval in $u$, independent of $R$, makes the
unchanged extracted twelve-root principal Gram positive on all the
above source boxes merely by decreasing $\eps_R$.
This obstruction disappears in the complete parent's $w=0$ slice.
\end{lemma}
\begin{proof}
Fix $z=R>1$, $b=1/(2R)$, $h=1/2$, $\sigs=t$, $x=t^2$, and let
$t\downarrow0$. All native gaps are positive for sufficiently small
$t$; in particular
\[
 D/t^2\to R/2,\quad \beta\to1/2,\quad t\lambda\to1,
 \quad d_G\to2R,\quad\rho\to1,\quad
 \nu\to R+1/2,\quad K_A\to R.
\]
Thus $c=0$ eventually. At $w=0$ the free-root principal block on
$(N_x,\Gamma_*)$ tends to
\[
 \begin{pmatrix}2&-u\sqrt{2R}\\-u\sqrt{2R}&1\end{pmatrix},
 \qquad \det\longrightarrow2(1-Ru^2).
\]
For any fixed $u>0$, choose $R>u^{-2}$ and then $t$ small enough to
enter any proposed positive source box. This proves the restricted
principal-Gram assertion.

It is essential that this is not an obstruction to the complete
certificate. At $w=0$, the complete parent is exactly
\begin{align*}
 F_w|_{w=0}={}&4m(W+uY)^2+4C(uW-Y)^2
       +4(A+\alpha)(1+u^2)Y^2\\
 &+2(1+u^2)(X^2+Y^2)
       +4D(uX-Z)^2+4\beta(X+uZ)^2 .
\end{align*}
This is a separated square identity for all $u$ and $z$; its testing
features may be chosen from $W,X,Y,Z,uW,uX,uY,uZ$.
The retained forms preserve linear relations between the roots that
are lost when the principal Gram is tested on free root variables.
For example, for arbitrary $w$,
\[
 uX_e=uN_x-Hu^2\Gamma-(Huw/q)A_\theta.
\]
The final angular term is a genuine additional coupling. No uniformly
bounded full-$w$ frame follows just by adding a scalar to the old
$(N_x,\Gamma_*)$ block.
\end{proof}

For a precise further target, the complete parent is quadratic in $w$:
$F_w=F_0+wF_1+w^2F_2$, where $F_0$ is the square identity above and
\begin{align*}
 F_1={}&-8m(W+uY)(X+uZ)
 +2(4D+2)(uX-Z)(uW-Y)\\
 &-2(4\alpha+2)(1+u^2)YZ,\\
 F_2={}&4m(X+uZ)^2+(4C+4D+2)(uW-Y)^2
 +4A(1+u^2)Y^2\\
 &+4\alpha(1+u^2)Z^2+(4\beta+2)(X+uZ)^2.
\end{align*}
Both $F_0$ and $F_2$ are separated squares with no $z$ restriction.
An $R$-independent testing extension therefore requires a uniform
completion of precisely the displayed cross form $F_1$, keeping
the native source relations and the testing identities. The present
note does not supply it, nor prove it impossible.

\paragraph{Exact ancillary checks.}
The adjacent script checks the two full-parent identities, all seven
source-path limits, the limiting minor, and the coefficient sum and
$z$ degree of $F_x$. Those exact algebraic checks do not replace the
inequality arguments above. No numerical semidefinite solve is used.
\end{document}
